\chapter{Conclusiones}
\label{cap:conclusiones}

Este capítulo reúne las respuestas a las cuatro preguntas de investigación, de la RQ1 a la RQ4, y evalúa los objetivos y las contribuciones con los resultados de los Capítulos~\ref{cap:catalogo} a~\ref{cap:comparacion}.
Las ocho Datasheets del catálogo pasan sin errores el validador de formato de este trabajo, que comprueba los tipos de cada campo y las referencias entre bloques, pero ninguna supera todas las comprobaciones que se le hacen al ingerirla en PRIME4API.
Al escribirlas aparecieron 35 elementos no soportados, hechos del dominio que ACM2YAML v0.3 no puede expresar o que PRIME4API no consume, como los seis volúmenes contratados que Postmark vende en un mismo plan.

El capítulo se organiza del siguiente modo.
En primer lugar, la Sección~\ref{sec:conclusiones:rq} responde a cada pregunta de investigación y la Sección~\ref{sec:conclusiones:evaluacion} da el estado de cada objetivo y de cada contribución.
A continuación, la Sección~\ref{sec:conclusiones:amenazas} agrupa las amenazas a la validez en cuatro ejes, cada una con su mitigación cuando la hay.
Después, la Sección~\ref{sec:conclusiones:lecciones} recoge las lecciones del método agéntico: el riesgo que añade un agente de inteligencia artificial, los controles que lo acotan y lo que no se ha medido.
Por último, la Sección~\ref{sec:conclusiones:futuro} plantea el trabajo futuro y el repositorio del que parte.

\section{Respuestas a las preguntas de investigación}\label{sec:conclusiones:rq}

La Tabla~\ref{tab:conclusiones:rq} responde a las cuatro preguntas de investigación de la Sección~\ref{sec:introduccion:objetivos}.
Cada celda da el resultado principal, y las cifras de detalle están en el capítulo que la sostiene, el que va entre paréntesis detrás de la pregunta.

\begin{table}[tbp]
\centering
\small
\begin{tabular}{|>{\raggedright\arraybackslash}p{3.3cm}|p{10.8cm}|}
\hline
\textbf{Pregunta} & \textbf{Respuesta} \\
\hline
RQ1. Qué APIs de correo incluir y con qué criterios reproducibles (Cap.~\ref{cap:catalogo}) &
El conjunto debe cubrir cuatro formas de cobro: volumen mensual, pago por uso, producto añadido y límite diario.
Sus APIs deben diferir en cómo limitan el envío, por ejemplo con un \textit{rate} por \textit{endpoint} o con un caudal negociado con el proveedor.
Las ocho del catálogo cumplen esa condición.
Es reproducible el protocolo de selección (tres fuentes de inclusión, cinco criterios ordenados y una regla de reemplazo).
La consulta al modelo de lenguaje no lo es, porque repetirla no garantiza las mismas APIs, y queda registrada en el Anexo~\ref{anx:uso-ia}. \\
\hline
RQ2. Si ACM2YAML basta para la capacidad contractual y qué extensiones requiere (Cap.~\ref{cap:evolucion}) &
No basta: las ocho Datasheets dejan en el registro de procedencia, que ninguna herramienta lee, datos como los seis volúmenes contratados de Postmark.
Las extensiones son nueve propuestas.
Siete amplían lo que el esquema expresa, y dos corrigen formas que la guía de ACM2YAML admite y PRIME4API no lee: el límite vacío (P-04) y el precio \texttt{custom} del plan Enterprise de Resend (P-05). \\
\hline
RQ3. Cómo se relacionan los límites de uso y los \textit{rates}, y qué implica para el consumidor (Cap.~\ref{cap:limites}) &
En seis de las ocho APIs la cuota es la restricción activa, el límite que corta antes el envío.
Con los planes que recoge la Datasheet, en lo que el \textit{rate} admite durante la ventana de la cuota cabe entre 103,7 veces la cuota (MailerSend Starter) y 129.600.000 (Postmark Free).
Quien contrata una de esas seis puede dimensionar el volumen sin consultar el \textit{rate} si su cuenta está aprobada y sin restricciones.
En Amazon SES y Mailchimp Transactional no se puede decir que mande la cuota.
El estado de la cuenta cambia la capacidad en seis de las ocho, esas dos incluidas, y ninguna publica el umbral que dispara el cambio. \\
\hline
RQ4. Cómo comparar las APIs con requisitos funcionales y no funcionales a la vez (Cap.~\ref{cap:comparacion}) &
Un procedimiento instancia en cada API el catálogo de 32 preguntas de decisión y responde cada una con el instrumento de su clase: la Datasheet, PRIME4API o HARVEY4API.
Son respondibles por su clase 25, y 21 sin HARVEY4API, que no se ha ejecutado.
El coste, la cuota y el caudal se calculan sobre la Datasheet, y los \textit{webhooks} o la IP dedicada se leen a mano en la matriz de \textit{features}.
En el caso de los 100.000 correos al mes cumplen cuatro APIs, y la más barata es SendGrid Essentials 100K, a 34,95 USD.
Los campos de ACM2YAML v0.3, con un solo volumen por plan, dan 2,4 veces ese precio. \\
\hline
\end{tabular}
\caption[Respuestas a las preguntas de investigación]{Respuesta a cada pregunta de investigación y capítulo que la sostiene.}
\label{tab:conclusiones:rq}
\end{table}

\section{Objetivos y contribuciones}\label{sec:conclusiones:evaluacion}

La Tabla~\ref{tab:conclusiones:evaluacion} da el estado de cada objetivo específico y de la contribución que le corresponde, con los epígrafes de la memoria que lo sostienen.
OBJ-03 y OBJ-04 se justifican primero, porque el Capítulo~\ref{cap:introduccion} sitúa en ellos el componente de investigación.

\begin{table}[tbp]
\centering
\small
\setlength{\tabcolsep}{3pt}
\begin{tabular}{|>{\raggedright\arraybackslash}p{3.0cm}|c|c|>{\raggedright\arraybackslash}p{2.0cm}|>{\raggedright\arraybackslash}p{2.2cm}|>{\raggedright\arraybackslash}p{3.0cm}|}
\hline
\textbf{Objetivo} & \textbf{RQ} & \textbf{Contrib.} & \textbf{Estado del objetivo} & \textbf{Estado de la contribución} & \textbf{Epígrafes que lo sostienen} \\
\hline
OBJ-01. Selección justificada del dominio & RQ1 & C-01 & cumplido en parte & cumplida en parte & \ref{sec:propuesta:seleccion}, \ref{sec:propuesta:respuesta} \\
\hline
OBJ-02. Construcción del catálogo de Datasheets & RQ2 a RQ4 & C-01 & cumplido & cumplida en parte & \ref{sec:propuesta:trazabilidad}, \ref{sec:propuesta:prime4api}, \ref{sec:esquema:contrato} \\
\hline
OBJ-03. Evolución del esquema & RQ2 & C-02 & cumplido en parte & cumplida & \ref{sec:esquema:inventario}, \ref{sec:esquema:fuera}, \ref{sec:esquema:respuesta} \\
\hline
OBJ-04. Análisis de la capacidad & RQ3 & C-03 & cumplido & cumplida & \ref{sec:limites:geometrica}, \ref{sec:limites:catalogo}, \ref{sec:limites:respuesta} \\
\hline
OBJ-05. Método de comparación & RQ4 & C-04 & cumplido & cumplida & \ref{sec:comparacion:catalogo}, \ref{sec:comparacion:procedimiento}, \ref{sec:comparacion:respuesta} \\
\hline
OBJ-06. Contribución a SPHERE & ninguna & C-01 & cumplido & cumplida en parte & \ref{sec:propuesta:sphere} \\
\hline
\end{tabular}
\caption[Estado de cada objetivo específico y de la contribución que le corresponde]{Estado de cada objetivo específico y de la contribución que le corresponde. RQ: la pregunta de investigación a la que sirve el objetivo. C-01: catálogo de Datasheets de APIs de correo. C-02: evolución de ACM2YAML. C-03: caracterización de la relación entre límites y \textit{rates}. C-04: método de comparación sistemática. OBJ-02 sirve de medio a RQ2, RQ3 y RQ4, porque el catálogo es la condición para responderlas. C-01 aparece tres veces porque responde a OBJ-01, OBJ-02 y OBJ-06. OBJ-06 no deriva de ninguna RQ.}
\label{tab:conclusiones:evaluacion}
\end{table}

\textbf{Evolución del esquema (OBJ-03 y C-02).}
De los 35 elementos no soportados que reúne el inventario de las ocho APIs, 24 superan el criterio de promoción: aparecen en dos APIs o más, o una ejecución de PRIME4API muestra su efecto sobre la Datasheet (Sección~\ref{sec:esquema:inventario}).
Hay nueve propuestas de extensión.
Ocho responden a elementos promovidos, como P-01, los seis volúmenes contratados que Postmark vende dentro de un mismo plan.
La novena, P-05, el plan con precio que existe y no es público, sale del fallo del motor ante el precio \texttt{custom} del plan Enterprise de Resend.

Los otros 16 elementos promovidos no tienen propuesta.
Doce solo tienen una forma esbozada que no se ha contrastado con la guía de ACM2YAML (Tabla~\ref{tab:sinforma}).
Los cuatro restantes quedan fuera de las propuestas por otras razones, como los \textit{add-ons} de IP dedicada, que el esquema deja al canal web (Sección~\ref{sec:esquema:fuera}).
OBJ-03 queda cumplido en parte, porque pide la extensión de cada carencia identificada y esos doce promovidos no la tienen fijada.

C-02 se cumple: cada una de las siete propuestas que amplían lo que ACM2YAML v0.3 puede expresar responde a un caso real que el esquema no representa, como el tope con que SendGrid limita el envío mensual a diez veces el volumen contratado, sobrecoste incluido.
Las otras dos quedan fuera del enunciado de C-02: P-04, la marca de valor no publicado junto al campo, y P-05.
ACM2YAML v0.3 admite esos valores, y es PRIME4API el que falla al leerlos.

\textbf{Análisis de la capacidad (OBJ-04 y C-03).}
La cuota es la restricción activa en Postmark, Resend, SendGrid, MailerSend, Brevo y Mailgun, seis de las ocho APIs (Sección~\ref{sec:limites:respuesta}).
La curva de capacidad $C(t)$ que dibuja PRIME4API (Figura~\ref{fig:capacidad}) da el máximo de correos que se pueden enviar desde el inicio de la ventana hasta cada instante sin pasar ningún límite.
Su codo es el instante en que deja de mandar el \textit{rate} y empieza a mandar la cuota, y con un destinatario por petición llega a los 83 minutos en Resend Pro y a los 41 en MailerSend Hobby.
Dos límites del catálogo no caben en esa curva: los dos contadores de Mailgun sobre la misma ventana y la franquicia sin ventana de Mailchimp Transactional, API en la que manda el caudal.

OBJ-04 y C-03 se cumplen sobre una base estrecha.
Solo en esas seis APIs la Datasheet da cifra al \textit{rate} y a la cuota de un mismo plan, y los valores medidos en MailerSend, Amazon SES y Mailchimp Transactional salen de cuentas sin aprobar o en pruebas, no del plan contratado.

\textbf{Método de comparación (OBJ-05 y C-04).}
El catálogo de preguntas reúne 32 preguntas de decisión, y el procedimiento que instancia cada una sobre cada API se ha aplicado a las ocho (Sección~\ref{sec:comparacion:procedimiento}).
OBJ-05 y C-04 se cumplen con tres salvedades.
La primera es que las cuatro preguntas que piden HARVEY4API cuentan como respondibles por la especificación de HARVEY4API que recoge Ruiz López \citep{ruizlopez2026tfm}, y solo se contrastaron con las respuestas que dio el modelo de lenguaje Claude Opus 5.5 al leer las Datasheets.

La segunda salvedad es que las diez respuestas de Mailchimp Transactional que resuelve una operación de PRIME4API están calculadas a mano.
El motor devuelve HTTP 500 en 14 llamadas por dos campos vacíos que el esquema admite: el tope de destinatarios y el periodo de la franquicia de la cuenta de prueba.
La tercera salvedad es que una de las tres preguntas funcionales, la retención de los registros de envío, pide un campo que ACM2YAML v0.3 no tiene.
Las otras dos se responden leyendo la matriz de \textit{features} por plan.

\textbf{Catálogo, selección y SPHERE (C-01 frente a OBJ-01, OBJ-02 y OBJ-06).}
OBJ-02, la construcción del catálogo, se cumple.
Los términos de consumo de cada proveedor quedan recogidos con su URL y su fecha de acceso, y cada valor de las ocho Datasheets lleva la cita literal que lo sostiene (Sección~\ref{sec:propuesta:trazabilidad}).
En Postmark no se archivaron las páginas que cita su Datasheet, y su versión se conserva solo como URL y fecha de acceso.

OBJ-01, la selección del dominio, se cumple en parte.
El protocolo parte de tres fuentes de inclusión: las APIs del trabajo de Noguerol Tirado \citep{noguerol2025tfm}, las ya modeladas en el ecosistema y una consulta a un modelo de lenguaje (Sección~\ref{sec:propuesta:fuentes}).
La primera no aportó ninguna API, porque se tuvo acceso a la memoria de Noguerol Tirado cuando la selección ya estaba cerrada.
La consulta presentaba Amazon SES como pago por uso puro, cuando ya publicaba planes, y repetirla no garantiza las mismas APIs (Sección~\ref{sec:propuesta:respuesta}).
Postmark, el piloto con el que se probó y se ajustó el método, no entra por ninguna de las tres fuentes.

OBJ-06, la contribución a SPHERE \citep{sphere}, se cumple.
Los ocho pricings en Pricing2YAML pasan el analizador de pricing4ts 0.10.4, la versión que usa la plataforma, y cada una de las 20 carencias del formato tiene su caso real.
Pricing2YAML no distingue un caudal de una cuota: en Resend Free, el caudal de 10 peticiones por segundo y la cuota de 3.000 correos al mes se escriben igual, como un límite renovable con su periodo, y la analítica de SPHERE los cuenta igual (Sección~\ref{sec:propuesta:sphere}).
La maqueta de la colección en SPHERE propone cómo mostrar los términos de consumo en la ficha de cada API.\footnote{La maqueta puede abrirse en \url{https://danvelcam.github.io/tfm-pricing-apis-correo/sphere/coleccion.html} (Sección~\ref{sec:propuesta:sphere:coleccion}).}

C-01 se cumple en parte por tres razones.
La primera es que la colección está cargada en SPHERE desde el 27 de octubre de 2026, pero es privada hasta el depósito.
La segunda es que ninguna Datasheet supera todas las comprobaciones que se le hacen al ingerirla en PRIME4API (Tabla~\ref{tab:verificacion}).
La tercera es que HARVEY4API, la otra herramienta del ecosistema que debería reutilizar las Datasheets, no se ha ejecutado.

En cinco de las ocho Datasheets, las comprobaciones fallan por un valor ausente o por un precio no público.
En Postmark queda vacío el tope de destinatarios, que el proveedor no pone, y en Mailchimp Transactional ese mismo tope y el periodo de la franquicia de la cuenta de prueba.
Fallan diez comprobaciones en cada una, y cinco en Resend, que no publica el \textit{rate} de su plan Enterprise.
Un precio que existe y no es público invalida trece comprobaciones en SendGrid, trece en Brevo y ocho en Resend.
MailerSend, Amazon SES y Mailgun solo fallan la comprobación de la curva de capacidad acumulada, con el mismo mensaje de error.
En MailerSend y Amazon SES el fallo es del motor, y en Mailgun no se ha aislado su causa.

\section{Amenazas a la validez}\label{sec:conclusiones:amenazas}

Esta sección agrupa por eje de validez las limitaciones que declaran los Capítulos~\ref{cap:catalogo} a~\ref{cap:comparacion}, cada una con su caso, su mitigación cuando la hay y lo que queda abierto.

\subsection{Validez de constructo}\label{sec:conclusiones:constructo}

La validez de constructo pregunta si lo medido es lo que se dice medir.
El sondeo de MailerSend midió el estado de la cuenta y no el plan: la cuenta sin aprobar admitió 10 peticiones por minuto frente a las 120 que el proveedor publica para una cuenta sana.
El sondeo de Amazon SES corrió en el entorno de pruebas (\textit{sandbox}) y el de Mailchimp Transactional en su cuenta de prueba.
Las tres Datasheets declaran el estado medido.
Medir en producción exige que el proveedor apruebe la cuenta, como en la salida del \textit{sandbox} de Amazon SES, o pagar un plan.

El \textit{rate} de Postmark y el de Mailgun salen de la cabecera de respuesta de una cuenta gratuita, sin provocar el código \texttt{429} con que la API rechaza el exceso, y la Datasheet los extiende a todos los planes.
En Postmark no lo provocaron 200 peticiones, y en Mailgun harían falta 80.000 en unos diez segundos.
Lo mitiga que las dos cifras vayan en ámbar, el color del semáforo para un valor interpretado y no literal.

SendGrid publica sus 10.000 peticiones por segundo en unas preguntas frecuentes y su envío no devuelve cabeceras de \textit{rate}, así que la cifra va en verde por ser literal, pero ningún sondeo la ha observado.
Queda abierta en las tres APIs una medición bajo carga, por ejemplo con el banco de pruebas de límites de uso de API-TaskForce, el grupo que publica la guía de ACM2YAML \citep{datasheetusagetesting}.

La curva de capacidad necesita la ventana de cada límite y lo que descuenta cada envío.
La ventana de la cuota, el periodo tras el que su contador vuelve a cero, está en rojo, el color del semáforo para un dato no publicado, en Mailgun, Brevo, Amazon SES y Mailchimp Transactional, y en ámbar en las otras cuatro.
Ninguna de las ocho declara la ventana del \textit{rate}, y los dos sondeos que la acotaron, en Brevo y Mailgun, no fijaron si es fija o deslizante.
Que cada destinatario descuente un correo de la cuota consta en Resend, SendGrid, Brevo y Amazon SES, y en Postmark, MailerSend, Mailgun y Mailchimp Transactional queda sin medir.
Esos huecos los marca el semáforo API a API, de modo que ninguno pasa por publicado sin serlo.

La cita literal de cada valor solo se contrasta de forma mecánica si se archivó la página, y de Postmark no se archivó ninguna de las que cita su Datasheet.
Lo mitigan la URL y la fecha de acceso de cada valor.

El ecosistema no tiene validador oficial de Datasheets, y la conformidad se comprueba con el validador de formato de este trabajo y con la ingestión en PRIME4API.
Los defectos de contrato, valores que ACM2YAML v0.3 admite y PRIME4API no sabe leer, los confirma un motor ajeno.
En cambio, las cuatro carencias que se promueven solo por recurrencia descansan en la lectura que este trabajo hace de la guía de ACM2YAML: los volúmenes contratados (P-01), el consumo de un lote de mensajes con varios destinatarios (P-02), el tope de una petición en otra unidad (P-03) y el estado de la suscripción (P-06).

P-06, el estado de la suscripción, muestra el riesgo de esa lectura.
La guía documenta alias de estado, bloques que dentro de un \textit{endpoint} dan a un estado de la cuenta sus propias cifras, con MailerSend como ejemplo, y la Datasheet de MailerSend añade uno para la cuenta sin aprobar.
En un alias, la guía solo documenta el \textit{rate} y la cuota, pero el código de referencia también lee un tope por petición.
El tope de dos destinatarios por petición medido en esa cuenta cabe, por tanto, en el esquema, y lo que falta es documentarlo en la guía.
P-06 sigue siendo una extensión porque el esquema no permite declarar el disparador del estado ni declarar el estado una sola vez para toda la Datasheet.

A cada pregunta se le asigna el instrumento que la responde, sea la Datasheet, PRIME4API o HARVEY4API, sabiendo ya lo que ofrece cada uno, y eso puede inflar la proporción de preguntas respondibles.
Por eso el 78~\% de preguntas respondibles, 25 de las 32 del catálogo de preguntas, no se compara de forma directa con el 54~\% que Cavero et al.\ obtienen sobre 56 plantillas con herramientas que ya existían \citep{cavero2026pricingintelligence} (Sección~\ref{sec:comparacion:catalogo}).
Lo mitiga que las preguntas sigan la gramática de plantillas de Cavero et al., de modo que pueden incorporarse a su banco de preguntas y medirse con sus herramientas.

El catálogo de preguntas se cerró en 32 sin un criterio de saturación fijado de antemano, aunque al pasar de tres APIs a ocho no ganó ninguna.

\subsection{Validez interna}\label{sec:conclusiones:interna}

La validez interna pregunta si los resultados dependen de las Datasheets o del montaje con que se obtienen.
El script que verifica las Datasheets contra PRIME4API clona el motor sin fijar la revisión, y un cambio del motor puede mover los errores de la Tabla~\ref{tab:verificacion} sin tocar la Datasheet.
Seis de las ocho ejecuciones registran la revisión \texttt{8e6011c} del repositorio de PRIME4API y pueden repetirse contra ella, y las de Brevo y Amazon SES no la registran.

El método de sondeo tiene tres puntos débiles.
En MailerSend, el \textit{endpoint} de cuota servía el saldo del día anterior once horas después de la reposición (Sección~\ref{sec:limites:catalogo}), así que el saldo que registra el sondeo al terminar puede ser el de la víspera.
En Brevo, el sondeo lee la cifra en el error con que el proveedor rechaza un destinatario imposible de entregar, pero con la cabecera de modo \textit{sandbox}, que solo valida el formato, el envío devolvió \texttt{201}.
Lo cubren un sondeo previo que comprueba que el proveedor rechaza la petición y un tope de envíos que cuenta como real cada petición aceptada.
En Mailchimp Transactional, desde once destinatarios o con un adjunto de 5~MB, el envío queda en cola, y por encima de ese umbral el sondeo no mide.

PRIME4API corrió en local con la misma imagen de contenedor que despliega SPHERE, y no en el despliegue público que cita Gavira Sánchez \citep{gavira2026tfm}, así que sus respuestas no se han contrastado contra ese servicio.
Las diez respuestas de Mailchimp Transactional que se calcularon a mano por el HTTP 500 del motor (Sección~\ref{sec:conclusiones:evaluacion}) no tienen confirmación de PRIME4API.

Los valores modelados dependen también de qué fuente se toma cuando dos discrepan.
La regla que da preferencia a la página de precios vale para las 37 divergencias, de las 82 registradas, en que esa página es uno de los lados.
Aun en ellas deja sin resolver las siete en que la página se contradice a sí misma, cuyo valor queda en ámbar.
Las 45 restantes enfrentan fuentes distintas de la página de precios: en la mayoría, dos fuentes del mismo proveedor, como su referencia de la API y su centro de ayuda, y en las demás, una de ellas frente a un sondeo o a un tercero.
Las ordenan otras reglas de modelado, como la que da preferencia al valor medido.
Dieciséis de esas 45 siguen abiertas, con las dos fuentes en el registro de procedencia.

Las Datasheets y el recuento de recurrencia los construye un agente de inteligencia artificial, que podría dar por leído un dato que la fuente no contiene.
Lo mitiga que el agente que construye no sea el que verifica, y queda abierto medir cuántos valores corrigió esa revisión (Sección~\ref{sec:conclusiones:lecciones}).

\subsection{Validez externa}\label{sec:conclusiones:externa}

La validez externa pregunta hasta dónde valen los resultados fuera del catálogo y de su ventana de toma de datos.
La respuesta negativa a la RQ2, si ACM2YAML v0.3 basta para representar las APIs de correo, no depende de la muestra, porque basta un caso que v0.3 no represente, como los seis volúmenes contratados de Postmark.

Sí depende de la muestra qué propuestas se promueven al esquema.
El tope de una petición medido en una unidad distinta de la suya, como los 50~MB de carga de \texttt{/email/bulk} en Postmark (P-03), aparece en las ocho APIs.
En cambio, tres propuestas descansan en una sola API: el tope del sobrecoste de SendGrid (P-07), el cupo de destinatarios de la cuenta sin aprobar de MailerSend (P-08) y la cola de Brevo por encima de la cuota (P-09).
Se promueven porque PRIME4API da en ellas un resultado que el proveedor contradice.
Al pasar de cuatro APIs a ocho ninguna propuesta perdió su apoyo.
Queda abierto repetir el inventario en otro dominio, comprobando cada propuesta con el caso que la motiva.

La marca de valor no publicado (P-04) y el plan con precio no público (P-05) se han medido contra un único motor, PRIME4API.
Claude Opus 5.5, con el que se contrastaron las preguntas de HARVEY4API, lee el YAML como texto y no cuenta como segundo motor.

En la relación entre límites de uso y \textit{rates} (RQ3), que la cuota sea la restricción activa se explica por las unidades, y por eso cabe esperarlo fuera del catálogo: frente a una cuota mensual o diaria, un \textit{rate} por segundo o por minuto, sostenido durante la ventana de la cuota, suma mucho más volumen.
Lo que sí depende de la muestra son las excepciones, como el \textit{rate} reducido de la cuenta sin aprobar de MailerSend (Sección~\ref{sec:conclusiones:constructo}) o la discrepancia de Mailgun.
Mailgun publica 1.500 peticiones por segundo por IP de cliente, y la cabecera del envío devuelve un límite de 80.000 peticiones en una ventana medida de unos diez segundos.
Que el \textit{rate} del envío individual no cambie con el plan vale en Resend, SendGrid y MailerSend, y Brevo lo contradice.

Las cuatro formas de cobro que el catálogo cubre (volumen mensual, pago por uso, producto añadido y límite diario) salen del primer criterio de selección, cubrir formas de cobro distintas (Sección~\ref{sec:propuesta:criterios}), y no de un censo del sector.
Entre las APIs de reserva de la selección, Zoho ZeptoMail cobra con créditos prepagados, una variante que el catálogo no incluye.

El catálogo modela solo el producto transaccional, el de los mensajes individuales que suele desencadenar el propio usuario, y sus conclusiones no se trasladan al de \textit{marketing}, que Resend factura por contactos y no por correos.

La oferta cambia con el tiempo: el repositorio de ACM2YAML conserva una Datasheet de MailerSend con precios que difieren de los publicados cinco meses después \citep{apisconsumptionmodel}.
Cada valor lleva su fecha de acceso, y falta un procedimiento que diga cuándo reverificar.

\subsection{Validez de conclusión}\label{sec:conclusiones:conclusion}

La validez de conclusión pregunta si los recuentos que sostienen las conclusiones están bien hechos.
La Tabla~\ref{tab:recurrencia} deja 31 combinaciones de elemento y API sin verificar y una en suspenso, las dos cuotas de Mailgun que ninguna página dice si se aplican a la vez.
Una aparición sin verificar no se cuenta, así que verificarla solo puede añadir APIs y cada recuento es un mínimo.
Bastaría confirmar una aparición más en cada uno para que pasaran a recurrentes tres elementos no promovidos: la región como unidad de gobierno, el tope fraccionario y los dos contadores sobre una misma ventana.
Ninguna promoción se retiraría\footnote{El detalle por fila está en \texttt{03\_datos/\allowbreak RECURRENCIA\_\allowbreak 2026-09-23.md}, en el repositorio del trabajo.}.

En Postmark, una respuesta «no publicado» cuenta como respondida en la pregunta del \textit{rate} máximo y no en la del límite diario ni en la de si el \textit{rate} sube con el plan.
Aplicar el mismo criterio a las tres movería una o dos de sus 18 respuestas.
La cobertura cuenta igual el ámbar que el verde, y el ámbar lo llevan 13 de las 18 respuestas de Postmark y 12 de las 20 de Mailgun, frente a 1 de 22 en Resend.

El mes se toma de 30 días.
Con 28, el menor cociente entre \textit{rate} y cuota de una cuenta sana, el de MailerSend Professional 2M, un volumen que la Datasheet no recoge, baja de 2,59 a 2,42.
Ninguna restricción activa cambia con la longitud del mes.

\section{Lecciones del método agéntico}
\label{sec:conclusiones:lecciones}

El Anexo~\ref{anx:uso-ia} declara qué tareas hicieron los agentes de inteligencia artificial y cuáles el autor.

Este trabajo separa a quien construye una Datasheet de quien la verifica.
El riesgo propio de un agente es dar por publicado un dato que la fuente vigente no contiene.
Le ocurrió a Claude Opus 5.5 en la prueba de la Sección~\ref{sec:comparacion:harvey}, que le planteó preguntas de comparación con las Datasheets en el prompt y sin ellas.
Sin ellas, atribuyó a Amazon SES en producción 50.000 correos al día y 14 por segundo.
La documentación vigente no da ninguna de las dos: la primera solo sale en una entrada del blog de AWS de 2014, y la segunda no está en ninguna fuente de su Datasheet.
El mismo error en un agente que construye una Datasheet pasaría al catálogo.

En este trabajo lo acotan dos controles.
El primero es que otro agente recibe la Datasheet terminada y la revisa campo a campo con las fuentes a la vista, en la fase de revisión del ciclo de la Sección~\ref{sec:propuesta:ciclo}.
El segundo es el verificador de citas, una comprobación automática que contrasta la cita literal de cada valor con la copia archivada de su página, y no se aplica a Postmark, cuyas páginas no se archivaron.
No se ha contado cuántos valores de las ocho Datasheets corrigieron esos dos controles, de modo que su eficacia no está medida.

El libro de criterios es el fichero de reglas común a las ocho Datasheets.
Cada duda de modelado que no cubría se decidió una sola vez y entró en él como regla nueva, de modo que no se volvió a plantear en otra API.
Así, los límites de recursos, como el número de dominios o de \textit{webhooks} de una cuenta, no se modelan como cuotas en ninguna de las ocho.

En la prueba de la Sección~\ref{sec:comparacion:harvey}, con las Datasheets en el prompt, Claude Opus 5.5 leyó bien los campos que pedían las cuatro preguntas, que tocan cinco de las ocho APIs.
Citó el campo de cada cifra, separó el valor publicado del interpretado y llegó a la misma conclusión que la referencia calculada a mano, la misma a la que llegó sin ellas.
La prueba comprueba la lectura y no que las Datasheets sean correctas: como se construyeron con la misma familia de modelos, un valor mal leído al construirlas pasaría a la respuesta sin que la prueba lo detectase.
Por eso un valor que solo ha deducido un agente se queda en ámbar, el color de lo interpretado, mientras no lo respalde un valor publicado o medido.

La reducción de horas del autor no se puede atribuir a los agentes.
Las horas suman \dato{total.real} frente a las \dato{total.prev} previstas, pero el tiempo de los agentes no se registró y no hay una medición del mismo trabajo sin ellos (Sección~\ref{sec:metodologia:desviaciones}).
Medir ese efecto exigiría registrar ese tiempo desde el inicio y construir sin agentes una muestra de Datasheets.

\section{Trabajo futuro}\label{sec:conclusiones:futuro}

Sobre ACM2YAML quedan tres líneas de trabajo.
La primera es distinguir el límite publicado del aplicado, el elemento F-04, que reúne los límites que el proveedor declara y que la respuesta a la petición no refleja (Sección~\ref{sec:esquema:fuera}).
Es el caso de Mailchimp Transactional y de Amazon SES, que aceptaron llamadas por encima del tope que declaran.
Queda por decidir qué cifra va a la Datasheet y cómo se declara la fuente de cada una.
La segunda es fijar las doce formas esbozadas de la Tabla~\ref{tab:sinforma}, de una fecha de verificación para cada afirmación del registro de procedencia (F-09) al canal que ofrece cada operación (F-22), y contrastarlas con la guía de ACM2YAML.

La tercera es llevar como \textit{issues} las nueve propuestas de extensión de la Tabla~\ref{tab:propuestas}, cada una con su caso y el número de APIs en que aparece, al repositorio que tiene que cambiar: el de la guía si cambia el esquema y el de PRIME4API si el arreglo está en el motor (Sección~\ref{sec:esquema:metodo}).
De esas nueve, las cinco que van de los volúmenes contratados al plan con precio no público (P-01 a P-05) no figuran en ningún \textit{issue} de los repositorios de la guía, de PRIME4API ni de SPHERE (Sección~\ref{sec:esquema:fuera}).
Las cuatro restantes, del estado de la suscripción al comportamiento de la cuota al superarse (P-06 a P-09), aún deben contrastarse con los \textit{issues} abiertos antes de enviarse.
Las Datasheets del catálogo irían además al repositorio de la guía, en una \textit{pull request}, junto a las que ya publica \citep{apisconsumptionmodel}.

En las herramientas de consulta queda ejecutar con HARVEY4API las cuatro preguntas de clase H, las que combinan varias Datasheets u operaciones de PRIME4API, y comparar sus respuestas con las de Claude Opus 5.5 y con el cálculo a mano (Sección~\ref{sec:comparacion:harvey}).
También queda vincular la Datasheet del canal API con el iPricing, el metamodelo con el que SPHERE representa los planes y las \textit{features} de la página de precios \citep{garciafernandez2024pricing4saas}.
Con él, la pregunta de qué API da IP dedicada al menor precio para 200.000 correos al mes, que en la Sección~\ref{sec:comparacion:funcional} se calcula a mano, podría responderse con una consulta a SPHERE.

En SPHERE, la primera línea es implementar la sección de términos de consumo que la Sección~\ref{sec:propuesta:sphere:terminos} deja diseñada.
Lo publicado y lo medido entrarían por una ingesta que guarda la fecha de cada página, en un historial separado de las versiones de \textit{pricing} que SPHERE ya mantiene.
Ese diseño separa los límites propios de una interfaz, como las 10 conexiones SMTP concurrentes por IP que recomienda Postmark, de los compartidos entre varias.
Un límite compartido es la cuota diaria de MailerSend, que cuenta las peticiones de la API REST y las del relé SMTP.

Al diseño le falta una tercera clase, la de los límites globales, que valen para el servicio entero y no para una interfaz concreta.
Esa clase exige un nivel por encima de la interfaz.

La segunda línea en SPHERE es comparar las vías de contratación de un mismo servicio.
En RapidAPI, un mercado de APIs, hay fichas del servicio de envío de cuatro de las ocho APIs: SendGrid, Postmark, Brevo y Mailgun\footnote{Búsqueda sin cuenta en el buscador de RapidAPI, consultada el 23 de octubre de 2026.}.
Comparar el precio y los límites de SendGrid contratado en RapidAPI con los de su página de precios pide un \textit{pricing} por vía de contratación y el árbol de proveedor, servicio e interfaz de la maqueta de la colección para SPHERE (Sección~\ref{sec:propuesta:sphere:coleccion}).
Queda por aclarar si esas fichas, que publica un tercero, exigen una cuenta en el proveedor.

Las últimas líneas son \textit{skills}, instrucciones que sigue un agente basado en un modelo de lenguaje para leer los términos de consumo, y van de la más próxima a este trabajo a la más lejana.
Una agruparía las interfaces que comparten un límite, como la cuota diaria de MailerSend, y su salida se contrastaría con la página antes de entrar en la Datasheet.
Otra listaría como candidatas las contradicciones entre páginas de un mismo proveedor, como la guía de envío por lotes de Brevo, que admite 6.000 llamadas por hora al \textit{endpoint} de envío, 600 veces menos que los 3.600.000 por hora de su categoría General (Sección~\ref{sec:limites:publicado}).
Una tercera contaría los límites de cada versión fechada de una página, como las capturas de la página de límites de MailerSend.

La cuarta traduciría a cambios del dominio el diff sintáctico con el que SPHERE compara dos versiones de un pricing, que marca qué líneas difieren entre ellas \citep{sphere}.
La más lejana, porque depende de las otras, es publicarlas en un \textit{marketplace} que registre su uso real y conectarlas con SPHERE por el \textit{Model Context Protocol}, el mismo con el que HARVEY4API publica las operaciones de PRIME4API \citep{mcp}.

\phantomsection\label{sec:conclusiones:replicacion}Cualquiera de estas líneas parte del repositorio de datos del trabajo, \url{https://github.com/danvelcam/tfm-pricing-apis-correo}.
Reúne las ocho Datasheets con sus fuentes y evidencias junto a las herramientas del trabajo.
No incluye las credenciales de las APIs ni PRIME4API, y de las copias archivadas de las páginas de los proveedores publica solo la huella sha256.

Para comprobar la instalación, la herramienta \path{05_salida/_herramientas/verificar_prime4api.py} se ejecuta con \texttt{python3} sobre \path{03_datos/apis/postmark/datasheet.yaml}.
La primera vez clona PRIME4API en \path{~/.cache/atf/prime4api} sin fijar la revisión, así que antes hay que situar esa copia en la revisión \texttt{8e6011c} (Sección~\ref{sec:conclusiones:interna}).
El resultado debe coincidir con la fila de Postmark de la Tabla~\ref{tab:verificacion}: 11 comprobaciones superadas y 10 fallidas por valor ausente.

